% Generated semantic source for AI Almanach v3.
% Edit v3 directly after initial generation; do not overwrite
% editorial changes by rerunning the converter casually.

% ==== Multi-Agent Systems =============================================================================
\chapter{Multi-Agent Systems}

\chapterlead{How can autonomous components coordinate without losing control?}{Architect}{Communicate}{Coordinate}{Verify}{multi-agent-systems}

% -------------------------------------------------------
\section{Agent Technology}

\subsection{Concepts of Agents and Multi-Agent Systems}

\begin{v3topic}{What is an Agent?}
    An \textbf{agent} is a computer system situated in some environment that is
    capable of \emph{autonomous action} in that environment in order to meet its
    design objectives.\textsuperscript{\cite[][p.~28]{wooldridgeIntroductionMultiAgentSystems2009}}

    Four key properties (RAPS):
    \begin{description}[leftmargin=0.6cm, labelindent=0pt]
        \item[\textbf{Reactivity}] Perceives its environment and responds in a timely fashion
        \item[\textbf{Autonomy}] Operates without direct human intervention; controls its own actions
        \item[\textbf{Proactivity}] Takes initiative; pursues goals, not just reactions
        \item[\textbf{Social ability}] Interacts with other agents (and humans) via some communication language
    \end{description}

\end{v3topic}
\begin{v3topic}{Multi-Agent System (MAS)}
    A \textbf{Multi-Agent System} (MAS) is a collection of agents that perceive
    and act in a shared environment. No single agent has full knowledge or control
    of the global
    system.\textsuperscript{\cite[][p.~2]{weissMultiagentSystems2013}}

    Key structural concepts:
    \begin{itemize}
        \item \textbf{Environment}: shared space agents act in (physical or virtual)
        \item \textbf{Interaction}: coordination, cooperation, competition
        \item \textbf{Organisation}: roles, norms, institutions governing agent behaviour
        \item \textbf{Emergence}: global behaviour arises from local interactions
    \end{itemize}

\end{v3topic}
\begin{v3topic}{Environment Types}
    Russell \& Norvig classify environments along five dimensions that determine agent design
    complexity:\textsuperscript{\cite[][p.~47]{wooldridgeIntroductionMultiAgentSystems2009}}
    \begin{center}
    \begin{tblr}{
        colspec={X[l]X[l]X[l]},
        hlines, vlines,
        row{1}={font=\bfseries\small, bg=LimeGreen!20},
        row{2-Z}={font=\small},
    }
        Dimension & Easier for agent & Harder for agent \\
        Observability & Fully observable & Partially observable \\
        Determinism & Deterministic (action outcome known) & Stochastic (uncertain outcomes) \\
        Episode structure & Episodic (no past dependence) & Sequential (history matters) \\
        Dynamics & Static (world waits) & Dynamic (world changes while deciding) \\
        Action/state space & Discrete & Continuous \\
    \end{tblr}
\captionof{table}{Concepts of Agents and Multi-Agent Systems: Environment Types.}
\label{tab:ch15-multi-agent-systems-inline-01}
    \end{center}

\end{v3topic}

\subsection{Agent Applications}

\begin{v3topic}{Application Domains}
    MAS have been successfully deployed across many real-world
    domains:\textsuperscript{\cite[][p.~6]{weissMultiagentSystems2013}}
    \begin{description}[leftmargin=0.6cm, labelindent=0pt]
        \item[\textbf{Autonomous vehicles}] Fleets coordinate merging, intersection clearing,
              platooning; each vehicle is an agent with perception--action loops
        \item[\textbf{Electronic trading}] Automated market-making, algorithmic portfolio agents
              competing and clearing on exchanges in milliseconds
        \item[\textbf{Game AI}] Non-player character squads, real-time strategy games (StarCraft),
              cooperative/adversarial game-playing
        \item[\textbf{Smart grid}] Demand-response agents balance load; prosumer agents trade
              surplus renewable energy peer-to-peer
        \item[\textbf{Manufacturing}] Scheduling agents allocate machines and workers; holonic
              agents organise self-reconfiguring production cells
        \item[\textbf{Logistics}] Route planning, warehouse robot swarms, container port scheduling
    \end{description}

\end{v3topic}
\begin{v3topic}{Why MAS?}
    MAS are the natural architecture when:
    \begin{itemize}
        \item The problem is \textbf{too large} for a single centralised solution
        \item Sub-systems are \textbf{geographically distributed}
        \item Multiple \textbf{autonomous stakeholders} must cooperate or compete
        \item \textbf{Robustness} to failures is required (no single point of failure)
        \item \textbf{Scalability} matters: adding new agents is incremental
    \end{itemize}
    \textsuperscript{\cite[][p.~3]{wooldridgeIntroductionMultiAgentSystems2009}}

\end{v3topic}

\subsection{Agent-Oriented Design}

\begin{v3topic}{BDI Architecture}
    The \textbf{Beliefs--Desires--Intentions} (BDI) model is the dominant
    deliberative agent architecture, inspired by folk psychology of rational
    action:\textsuperscript{\cite[][p.~73]{wooldridgeIntroductionMultiAgentSystems2009}}
    \begin{description}[leftmargin=0.6cm, labelindent=0pt]
        \item[\textbf{Beliefs}] Agent's knowledge about itself and the world (possibly incomplete/incorrect)
        \item[\textbf{Desires}] Motivational states; goals the agent would like to achieve
        \item[\textbf{Intentions}] Commitments to pursue: current goals the agent has decided to act upon
    \end{description}
    The agent \emph{deliberates} to select intentions from desires, then
    \emph{plans} to find action sequences that achieve them.
    Practical implementations include \textbf{AgentSpeak} and \textbf{Jason}.

\end{v3topic}
\begin{v3topic}{FIPA Standards}
    \textbf{FIPA} (Foundation for Intelligent Physical Agents) defined
    interoperability standards still widely used:\textsuperscript{\cite[][Ch.~2]{bellifeministeDevelopingMultiAgentSystems2007}}
    \begin{itemize}
        \item \textbf{FIPA-ACL}: Agent Communication Language with performatives
              (inform, request, propose, agree, refuse, \ldots)
        \item \textbf{Abstract Architecture}: platform services — Agent Management
              System (AMS), Directory Facilitator (DF), Message Transport System (MTS)
        \item \textbf{Interaction Protocols}: standardised multi-step conversation patterns
              (FIPA Request, Contract Net, Auctions)
        \item \textbf{JADE}: open-source Java middleware fully compliant with FIPA
    \end{itemize}

\end{v3topic}
\begin{v3topic}{Agent-Oriented Methodologies}
    Several methodologies guide the \emph{analysis and design} of MAS, analogous to
    UML for object-oriented
    systems:\textsuperscript{\cite[][p.~9]{weissMultiagentSystems2013}}
    \begin{center}
    \begin{tblr}{
        colspec={X[l]X[l]X[l]},
        hlines, vlines,
        row{1}={font=\bfseries\small, bg=LimeGreen!20},
        row{2-Z}={font=\small},
    }
        Methodology & Key Abstraction & Outputs \\
        \textbf{Gaia} & Roles \& protocols & Role model, interaction model \\
        \textbf{Prometheus} & Goals, percepts, actions & System overview, detailed design \\
        \textbf{MaSE} (Multi-Agent SE) & Use cases $\to$ agents & Deployment, sequence diagrams \\
        \textbf{Tropos} & Social dependencies & Goal model, agent capabilities \\
    \end{tblr}
\captionof{table}{Agent-Oriented Design: Agent-Oriented Methodologies.}
\label{tab:ch15-multi-agent-systems-inline-02}
    \end{center}
    All start from a \emph{societal} view (roles, norms, dependencies) and
    progressively refine down to individual agent implementation.

\end{v3topic}

% -------------------------------------------------------
\section{Types of Intelligent Agents}

\subsection{Reasoning (Deliberative) Agents}

\begin{v3topic}{Symbolic Reasoning Agents}
    \textbf{Deliberative agents} maintain an explicit internal symbolic model
    of the world and use logical \emph{inference} to decide on
    actions.\textsuperscript{\cite[][p.~58]{wooldridgeIntroductionMultiAgentSystems2009}}

    Core components:
    \begin{itemize}
        \item \textbf{World model}: first-order logic or state-space representation
        \item \textbf{Goal representation}: explicit target states or utility functions
        \item \textbf{Planner}: finds action sequences leading to goal states
        \item \textbf{Deliberation loop}: perceive $\to$ update beliefs $\to$ select goal $\to$ plan $\to$ execute
    \end{itemize}
    Strength: principled, explainable.
    Weakness: computational cost (symbol grounding, planning is PSPACE-hard); brittle under incomplete information.

\end{v3topic}
\begin{v3topic}{Planning: STRIPS \& PDDL}
    \textbf{STRIPS} (Stanford Research Institute Problem Solver) is the classical
    planning formalism: actions described by \emph{preconditions} and
    \emph{effects}.\textsuperscript{\cite[][p.~62]{wooldridgeIntroductionMultiAgentSystems2009}}

    \begin{itemize}
        \item State = set of ground atoms (fluents)
        \item Action: \texttt{Pre} (must hold), \texttt{Add} (set true), \texttt{Del} (set false)
        \item Plan: sequence of actions that transitions initial state to goal state
    \end{itemize}
    \textbf{PDDL} (Planning Domain Definition Language) generalises STRIPS to
    typed objects, numeric fluents, durative actions, and temporal constraints.
    Used in international planning competitions (IPC).

\end{v3topic}

\subsection{Reactive Agents}

\begin{v3topic}{Subsumption Architecture (Brooks)}
    Rodney Brooks argued against explicit models: intelligence emerges from
    \emph{embodied} interaction with the physical
    world.\textsuperscript{\cite[][p.~66]{wooldridgeIntroductionMultiAgentSystems2009}}

    \textbf{Subsumption architecture}:
    \begin{itemize}
        \item Behaviour is decomposed into a \textbf{hierarchy of layers}
        \item Each layer is a simple \textbf{stimulus--response rule}: percept $\to$ action
        \item Higher layers \emph{subsume} (override) lower layers
        \item No internal world model, no planning
        \item Example layers (bottom-up): \emph{avoid obstacles} $\to$ \emph{wander} $\to$ \emph{explore}
    \end{itemize}
    Reactive agents are fast and robust but cannot reason about goals not
    immediately achievable by current percepts.

\end{v3topic}
\begin{v3topic}{Reactive vs.\ Deliberative}
    \begin{tblr}{
        colspec={X[l]X[l]},
        hlines, vlines,
        row{1}={font=\bfseries\small, bg=LimeGreen!20},
        row{2-Z}={font=\small},
    }
        Reactive & Deliberative \\
        No internal model & Explicit world model \\
        Stimulus--response rules & Planning \& reasoning \\
        Fast response & May be slow (planning cost) \\
        Robust to noise & Brittle if model wrong \\
        Limited to present percepts & Can pursue long-horizon goals \\
        Simple to implement & Complex engineering \\
    \end{tblr}
\captionof{table}{Reactive Agents: Reactive vs.\ Deliberative.}
\label{tab:ch15-multi-agent-systems-inline-03}
    \textsuperscript{\cite[][p.~67]{wooldridgeIntroductionMultiAgentSystems2009}}

\end{v3topic}

\subsection{Hybrid Agents}

\begin{v3topic}{Layered Hybrid Architectures}
    \textbf{Hybrid agents} combine a reactive layer for fast sensor-motor responses
    with a deliberative layer for goal reasoning and planning --- each layer
    operates at a different time
    scale.\textsuperscript{\cite[][p.~71]{wooldridgeIntroductionMultiAgentSystems2009}}

    \textbf{TouringMachines} (Ferguson, 1992): three concurrent layers
    connected by a control subsystem:
    \begin{enumerate}
        \item \textbf{Reactive}: immediate reflexes
        \item \textbf{Planning}: achieve current goals
        \item \textbf{Modelling}: predict and reason about other agents
    \end{enumerate}
    The control subsystem mediates which layer's action is actually executed.

\end{v3topic}
\begin{v3topic}{InteRRaP Architecture}
    \textbf{InteRRaP} (Müller, 1996) is a three-layer vertical control
    architecture:\textsuperscript{\cite[][p.~72]{wooldridgeIntroductionMultiAgentSystems2009}}

    \begin{enumerate}
        \item \textbf{Behaviour-based layer}: reactive, environment model
        \item \textbf{Local planning layer}: goal-directed planning, mental model of self
        \item \textbf{Cooperative planning layer}: social model of other agents; negotiation
    \end{enumerate}
    Control flows \emph{bottom-up} when lower layers fail; higher layers can
    activate lower-layer behaviours \emph{top-down}.
    This architecture is prototypical for warehouse and manufacturing agents.

\end{v3topic}

% -------------------------------------------------------
\section{Agent Communication}

\subsection{Ontologies}

\begin{v3topic}{Ontologies in MAS}
    Agents must share a common vocabulary to communicate meaningfully.
    An \textbf{ontology} formally specifies concepts, relationships, and
    constraints in a domain.\textsuperscript{\cite[][p.~198]{wooldridgeIntroductionMultiAgentSystems2009}}

    \begin{description}[leftmargin=0.6cm, labelindent=0pt]
        \item[\textbf{RDF}] Resource Description Framework: subject--predicate--object triples
        \item[\textbf{OWL}] Web Ontology Language: builds on RDF; supports classes, properties,
              axioms, description logic reasoning
        \item[\textbf{SPARQL}] Query language for RDF triple stores
    \end{description}
    Domain ontologies (e.g., SUMO, BioPortal medical ontologies) enable
    interoperability between heterogeneous agents: agents can translate
    between their local representations using shared ontology terms.

\end{v3topic}
\begin{v3topic}{Knowledge Representation Levels}
    Shared understanding in MAS requires grounding at multiple levels:
    \begin{enumerate}
        \item \textbf{Lexical}: same symbols used
        \item \textbf{Semantic}: same meaning attached to symbols
        \item \textbf{Pragmatic}: same inferences drawn from symbols
    \end{enumerate}
    Without semantic alignment, agents may use the same word (\texttt{Price})
    with conflicting definitions (inclusive/exclusive VAT).

    \textbf{Ontology alignment} and \textbf{ontology mapping} tools (e.g., PROMPT,
    GLUE) automatically match concepts across ontologies, enabling federation
    of heterogeneous agent knowledge
    bases.\textsuperscript{\cite[][p.~201]{weissMultiagentSystems2013}}

\end{v3topic}

\subsection{Communication Languages}

\begin{v3topic}{Speech Act Theory}
    Agent communication languages are grounded in Austin and Searle's
    \textbf{speech act theory}: utterances are \emph{actions} with
    intentional content.\textsuperscript{\cite[][p.~184]{wooldridgeIntroductionMultiAgentSystems2009}}

    A message has three layers:
    \begin{itemize}
        \item \textbf{Locution}: the literal content of the message
        \item \textbf{Illocution}: the speech act performed (\emph{performative}):
              \texttt{inform}, \texttt{request}, \texttt{propose}, \texttt{agree},
              \texttt{refuse}, \texttt{query-if}, \ldots
        \item \textbf{Perlocution}: the intended effect on the receiver's mental state
    \end{itemize}
    Both KQML and FIPA-ACL build on this model by wrapping content in a
    performative envelope.

\end{v3topic}
\begin{v3topic}{FIPA-ACL Message Structure}
    A \textbf{FIPA-ACL} message carries:\textsuperscript{\cite[][Ch.~3]{bellifeministeDevelopingMultiAgentSystems2007}}
    \begin{itemize}
        \item \texttt{performative}: illocutionary force (e.g., \texttt{INFORM})
        \item \texttt{sender} / \texttt{receiver}: agent identifiers
        \item \texttt{content}: domain-specific payload (string, fipa-sl expression, \ldots)
        \item \texttt{language}: content encoding language (e.g., FIPA-SL)
        \item \texttt{ontology}: shared vocabulary identifier
        \item \texttt{conversation-id}: ties messages into an interaction protocol
        \item \texttt{reply-with} / \texttt{in-reply-to}: threading
    \end{itemize}

\end{v3topic}
\begin{v3topic}{Interaction Protocols}
    Interaction protocols define the expected \emph{sequence of messages}
    between agents --- the agent-level analogue of network protocols.
    \textsuperscript{\cite[][Ch.~4]{bellifeministeDevelopingMultiAgentSystems2007}}
    \factvisual{
        \textbf{A performative has protocol meaning.}
        A request is not complete when it is sent. The receiver may refuse or
        agree, and an agreement must eventually resolve as information or
        failure.
        \begin{itemize}
            \item Conversation IDs correlate asynchronous replies.
            \item Timeouts and failure paths are part of the protocol.
            \item Payload validity does not guarantee semantic agreement.
        \end{itemize}
    }{
        \begin{tikzpicture}[x=1cm,y=0.72cm]
            \node[font=\small\sffamily\bfseries] at (0,4.8) {Requester};
            \node[font=\small\sffamily\bfseries] at (4.5,4.8) {Provider};
            \draw[AlmanachInk!45] (0,4.4) -- (0,0);
            \draw[AlmanachInk!45] (4.5,4.4) -- (4.5,0);
            \draw[almanach arrow] (0,3.9) -- node[above,font=\scriptsize]
                {REQUEST} (4.5,3.9);
            \draw[almanach arrow] (4.5,3.0) -- node[above,font=\scriptsize]
                {AGREE / REFUSE} (0,3.0);
            \draw[almanach arrow] (4.5,1.6) -- node[above,font=\scriptsize]
                {INFORM / FAILURE} (0,1.6);
            \node[draw=AlmanachOchre,dashed,rounded corners,
                  fit={(0,3.25)(4.5,1.25)},inner sep=5pt,
                  label={[font=\tiny\sffamily,text=AlmanachOchre]below:
                  one conversation}] {};
        \end{tikzpicture}
        \visualcaption{The FIPA Request interaction as a sequence: acceptance
          and task outcome are separate messages.}{fig:fipa-request-sequence}
    }
    \begin{center}
    \begin{tblr}{
        colspec={X[l]X[l]X[l]},
        hlines, vlines,
        row{1}={font=\bfseries\small, bg=LimeGreen!20},
        row{2-Z}={font=\small},
    }
        Protocol & Flow & Typical Use \\
        FIPA Request & REQUEST $\to$ AGREE/REFUSE $\to$ INFORM/FAILURE & Simple service invocation \\
        Contract Net & Announce $\to$ bids $\to$ award & Task allocation to best bidder \\
        Iterated Contract Net & Multiple rounds & Progressive negotiation \\
        English Auction & Ascending bids until timeout & Resource allocation \\
        Dutch Auction & Price decreasing until acceptance & Perishable goods allocation \\
        Brokering / Recruiting & Mediator finds capable agents & Service discovery \\
    \end{tblr}
\captionof{table}{Communication Languages: Interaction Protocols.}
\label{tab:ch15-multi-agent-systems-inline-04}
    \end{center}

\end{v3topic}

% -------------------------------------------------------
\section{Agent Cooperation}

\subsection{Distributed Problem Solving}

\begin{v3topic}{Task Decomposition and Sharing}
    \textbf{Distributed Problem Solving} (DPS): a problem too large for one
    agent is decomposed into sub-tasks distributed across cooperating agents.
    \textsuperscript{\cite[][p.~154]{weissMultiagentSystems2013}}

    Two fundamental patterns:
    \begin{description}[leftmargin=0.6cm, labelindent=0pt]
        \item[\textbf{Task sharing}] Work is divided among agents; each solves a sub-problem.
              Requires task allocation (Contract Net, market-based)
        \item[\textbf{Result sharing}] All agents partially observe the problem; they share
              partial results to produce a combined solution (sensor fusion, evidence aggregation)
    \end{description}
    Combination: agents solve different sub-tasks \emph{and} share intermediate
    results to correct each other.

\end{v3topic}
\begin{v3topic}{Blackboard Systems}
    A \textbf{blackboard system} is a shared data structure (the
    \emph{blackboard}) on which specialist agents post and read partial
    results.\textsuperscript{\cite[][p.~157]{weissMultiagentSystems2013}}

    Architecture:
    \begin{enumerate}
        \item \textbf{Blackboard}: hierarchically structured global memory
        \item \textbf{Knowledge Sources} (KS): specialist agents that read/write blackboard entries
        \item \textbf{Control}: scheduler selects which KS fires next
    \end{enumerate}
    The blackboard mediates all inter-agent communication --- agents are
    decoupled; they only know the shared data structure, not each other.
    Classic application: Hearsay-II speech understanding system (1970s).

\end{v3topic}

\subsection{Coordination Mechanisms}

\begin{v3topic}{Coordination and Inconsistency}
    In a MAS, agents may develop \textbf{conflicting plans} or
    \textbf{redundant work} without coordination.
    \textsuperscript{\cite[][p.~172]{wooldridgeIntroductionMultiAgentSystems2009}}

    Sources of inconsistency:
    \begin{itemize}
        \item \textbf{Resource conflicts}: multiple agents need the same resource simultaneously
        \item \textbf{Effect conflicts}: agent A's action undoes agent B's effect
        \item \textbf{Incomplete knowledge}: agents hold contradictory beliefs
    \end{itemize}
    Mechanisms to handle them:
    \begin{itemize}
        \item \textbf{Commitment protocols}: agents announce and respect published intentions
        \item \textbf{Partial Global Planning} (PGP): exchange local plans and merge globally
        \item \textbf{Synchronisation primitives}: lock-like coordination on shared resources
    \end{itemize}

\end{v3topic}
\begin{v3topic}{Asynchronous Teams}
    \textbf{Asynchronous Teams} (A-Teams, Talukdar et al.) are networks of agents
    that cooperate by \emph{improving a shared set of solutions} stored in a
    common memory.\textsuperscript{\cite{talukdarASYNCHRONOUSTEAMSCOOPERATION}}

    Key properties:
    \begin{itemize}
        \item Agents operate \textbf{asynchronously}: no synchronisation required
        \item Agents are \textbf{autonomous}: each applies its own heuristic to solutions
        \item Agents communicate \textbf{indirectly}: by writing improved solutions to shared memory
        \item Effective for combinatorial optimisation (routing, scheduling)
    \end{itemize}
    A-Teams combine the diversity of multi-heuristic search with the
    accumulated improvement of cooperative evolution.

\end{v3topic}

\subsection{Negotiation}

\begin{v3topic}{Bargaining and Monotonic Concession}
    \textbf{Negotiation} is a process by which agents reach mutually acceptable
    agreements on the division of tasks or
    resources.\textsuperscript{\cite[][p.~185]{wooldridgeIntroductionMultiAgentSystems2009}}

    \textbf{Monotonic Concession Protocol} (Zeuthen strategy):
    \begin{enumerate}
        \item Each agent opens with its preferred proposal
        \item At each round, at least one agent \emph{concedes}: moves to a proposal
              less preferred by itself but closer to the other agent's position
        \item No agent ever proposes an option worse than its previous proposal
        \item Agreement when proposals coincide, or one agent accepts
    \end{enumerate}
    Pareto-optimal outcomes are reached when agents concede efficiently.

\end{v3topic}
\begin{v3topic}{Argumentation-Based Negotiation}
    Beyond simple proposals, agents can use \textbf{argumentation} to justify
    positions and challenge counter-proposals:\textsuperscript{\cite[][p.~191]{wooldridgeIntroductionMultiAgentSystems2009}}

    \begin{description}[leftmargin=0.6cm, labelindent=0pt]
        \item[\textbf{Threats}]: ``If you do not concede on price, I will leave''
        \item[\textbf{Rewards}]: ``If you agree to this price, I offer future business''
        \item[\textbf{Appeal to precedent}]: ``We agreed to similar terms last time''
        \item[\textbf{Appeal to self-interest}]: explaining why agreement benefits receiver
    \end{description}
    Argumentation-based negotiation allows agents to \emph{change each other's beliefs},
    not just exchange proposals, leading to richer outcomes than pure bargaining.

\end{v3topic}

% -------------------------------------------------------
\section{Multi-Agent Decision-Making}

\subsection{Game Theory for Agents}

\begin{v3topic}{Normal-Form Games}
    A \textbf{normal-form game} is the basic model of strategic interaction:
    \textsuperscript{\cite[][Ch.~3]{shohamMultiagentSystemsAlgorithmic2009}}
    \[
        G = (N,\, A_1 \times \cdots \times A_n,\, u_1, \ldots, u_n)
    \]
    where $N$ is the set of agents, $A_i$ the action set of agent $i$, and
    $u_i: A_1 \times \cdots \times A_n \to \mathbb{R}$ the utility function.
    Agents choose actions \emph{simultaneously}; payoffs depend on the joint action.

    Classic examples:
    \begin{itemize}
        \item \textbf{Prisoner's Dilemma}: rational self-interest leads to a
              collectively suboptimal outcome
        \item \textbf{Coordination Game}: multiple equilibria; agents must align
        \item \textbf{Zero-sum}: one agent's gain is another's loss
    \end{itemize}

\end{v3topic}
\begin{v3topic}{Nash Equilibrium}
    A \textbf{Nash Equilibrium} is a joint strategy profile $\sigma^* =
    (\sigma_1^*, \ldots, \sigma_n^*)$ such that no agent can improve its payoff
    by unilaterally deviating:\textsuperscript{\cite[][p.~67]{shohamMultiagentSystemsAlgorithmic2009}}
    \begin{equation}
        u_i(\sigma_i^*,\, \sigma_{-i}^*) \;\geq\; u_i(\sigma_i,\, \sigma_{-i}^*)
        \quad \forall\, \sigma_i,\; \forall\, i \in N
    \end{equation}
    where $\sigma_{-i}^*$ denotes all agents' strategies except agent $i$.
    \begin{itemize}
        \item \textbf{Pure NE}: each agent plays one action with certainty
        \item \textbf{Mixed NE}: agents randomise over actions (Nash's theorem guarantees
              at least one mixed NE in every finite game)
        \item NE does not imply efficiency: Prisoner's Dilemma NE is Pareto-inferior
    \end{itemize}

\end{v3topic}
\begin{v3topic}{Dominant Strategies \& Pareto Optimality}
    Two key solution concepts complement Nash Equilibrium:\textsuperscript{\cite[][p.~64]{shohamMultiagentSystemsAlgorithmic2009}}

    \textbf{Dominant strategy}: action $a_i^*$ that is best for agent $i$ \emph{regardless}
    of what others do:
    \begin{equation}
        u_i(a_i^*, a_{-i}) \;\geq\; u_i(a_i, a_{-i})
        \quad \forall\, a_i \in A_i,\; \forall\, a_{-i}
    \end{equation}
    When every agent has a dominant strategy, those strategies form a
    \emph{dominant strategy equilibrium} (often unique and robust).

    \textbf{Pareto optimality}: outcome $o$ is Pareto optimal if no other outcome
    $o'$ makes \emph{every} agent at least as well off and at least one agent
    strictly better off.
    A Nash equilibrium need not be Pareto optimal (Prisoner's Dilemma);
    mechanism design seeks to align the two.

\end{v3topic}

\subsection{Social Choice and Coalitions}

\begin{v3topic}{Voting Mechanisms}
    \textbf{Social choice} aggregates individual preferences into a collective
    decision:\textsuperscript{\cite[][Ch.~9]{shohamMultiagentSystemsAlgorithmic2009}}
    \begin{description}[leftmargin=0.6cm, labelindent=0pt]
        \item[\textbf{Plurality}] Each agent votes for one option; most votes wins
        \item[\textbf{Borda count}] Each agent ranks all options; points awarded by rank
        \item[\textbf{Condorcet}] Option that beats all others in pairwise majority
        \item[\textbf{Approval voting}] Agents approve any number; most approvals wins
    \end{description}
    Arrow's Impossibility Theorem shows no deterministic voting rule simultaneously
    satisfies unanimity, independence of irrelevant alternatives, and non-dictatorship.

\end{v3topic}
\begin{v3topic}{Coalition Formation and Shapley Value}
    A \textbf{coalition} is a subset $C \subseteq N$ of agents that agree to
    coordinate. \textbf{Coalition formation} is the process of deciding which
    coalitions form and how to distribute the
    value.\textsuperscript{\cite[][Ch.~12]{shohamMultiagentSystemsAlgorithmic2009}}

    The \textbf{Shapley value} provides a fair allocation of coalition payoff
    $v(N)$ to agent $i$:
    \begin{equation}
        \phi_i(v) = \sum_{S \subseteq N \setminus \{i\}}
        \frac{|S|!\,(|N|-|S|-1)!}{|N|!}
        \bigl[v(S \cup \{i\}) - v(S)\bigr]
    \end{equation}
    This is the average marginal contribution of agent $i$ over all possible
    orderings of agents joining the coalition.
    Properties: efficiency ($\sum_i \phi_i = v(N)$), symmetry, dummy player, additivity.

\end{v3topic}

\subsection{Auction Mechanisms}

\begin{v3topic}{Auction Types}
    Auctions are \textbf{resource allocation mechanisms} with built-in incentive
    properties:\textsuperscript{\cite[][Ch.~11]{shohamMultiagentSystemsAlgorithmic2009}}
    \begin{description}[leftmargin=0.6cm, labelindent=0pt]
        \item[\textbf{English (ascending)}] Open cry; price rises; last bidder wins.
              Dominant strategy: bid up to true value
        \item[\textbf{Dutch (descending)}] Price falls from high; first bidder wins.
              Fast but requires strategic timing
        \item[\textbf{First-price sealed-bid}] Highest sealed bid wins at own bid.
              Bidders shade below true value
        \item[\textbf{Vickrey (second-price sealed-bid)}] Highest bid wins but pays
              \emph{second}-highest bid. \textbf{Dominant strategy}: truthful bidding
              (incentive-compatible)
    \end{description}

\end{v3topic}
\begin{v3topic}{Mechanism Design}
    \textbf{Mechanism design} (``reverse game theory'') asks: given a desired
    social outcome, what rules of the game achieve it as an
    equilibrium?\textsuperscript{\cite[][p.~280]{shohamMultiagentSystemsAlgorithmic2009}}

    Key desiderata:
    \begin{itemize}
        \item \textbf{Incentive compatibility}: truthful reporting is dominant strategy
        \item \textbf{Individual rationality}: agents voluntarily participate
        \item \textbf{Efficiency}: total welfare maximised
        \item \textbf{Budget balance}: no external subsidy required
    \end{itemize}
    The \textbf{Vickrey--Clarke--Groves} (VCG) mechanism achieves
    efficiency and incentive compatibility; combinatorial auctions use VCG
    for multi-item allocation.

\end{v3topic}

% -------------------------------------------------------
\section{Multi-Agent Reinforcement Learning}

\subsection{From Single-Agent to Multi-Agent RL}

\begin{v3topic}{The Non-Stationarity Problem}
    In single-agent RL the environment is stationary: the agent's policy
    does not change the transition or reward function.
    In a MAS, each agent's \emph{learning} changes the effective environment
    of every other agent, creating
    \textbf{non-stationarity}.\textsuperscript{\cite[][Ch.~17]{suttonReinforcementLearningIntroduction2018}}

    Consequence: the Markov property is violated from each agent's perspective.
    \begin{itemize}
        \item \textbf{Independent Q-Learners} (IQL): each agent runs Q-learning
              treating others as part of the environment; simple but convergence
              not guaranteed due to non-stationarity
        \item \textbf{Joint action learners}: track all agents' actions but suffer
              from exponential action space growth
    \end{itemize}

\end{v3topic}
\begin{v3topic}{Centralised Training, Decentralised Execution (CTDE)}
    \textbf{CTDE} is the dominant MARL paradigm:
    during \emph{training} a centralised critic has access to global state and
    all agents' actions; during \emph{execution} each agent acts using only
    its local observation.\textsuperscript{\cite[][p.~450]{weissMultiagentSystems2013}}

    Benefits:
    \begin{itemize}
        \item Centralised critic resolves non-stationarity at training time
        \item Decentralised execution is scalable and robust (no communication needed)
        \item Natural fit for cooperative tasks (factory robots, drone swarms)
    \end{itemize}
    Key challenge: \textbf{credit assignment} --- how much did each agent
    contribute to the shared reward?

\end{v3topic}

\subsection{Key MARL Algorithms}

\begin{v3topic}{Value Decomposition: QMIX}
    \textbf{QMIX} (Rashid et al., 2018) addresses cooperative MARL with a shared
    reward by learning individual agent Q-values $Q_i(o_i, a_i)$ that combine
    into a global $Q_{tot}$ via a \emph{mixing
    network}:\textsuperscript{\cite[][p.~452]{weissMultiagentSystems2013}}

    \begin{equation}
        Q_{tot}(\boldsymbol{\tau}, \mathbf{a}) = f_{\text{mix}}\bigl(Q_1(\tau_1, a_1), \ldots, Q_n(\tau_n, a_n);\, s\bigr)
    \end{equation}
    The mixing network has \emph{non-negative weights} (monotonicity constraint),
    ensuring the individual-global-max (IGM) principle:
    \[
        \arg\max_{\mathbf{a}} Q_{tot} = \bigl(\arg\max_{a_1} Q_1, \ldots, \arg\max_{a_n} Q_n\bigr)
    \]
    Each agent can act greedily without needing others' information.

\end{v3topic}
\begin{v3topic}{MADDPG and MAPPO}
    \textbf{MADDPG} (Lowe et al., 2017) extends DDPG to the multi-agent setting
    via CTDE: each agent has a centralised critic that receives all agents'
    observations and actions, but a decentralised actor that uses only local
    observation.\textsuperscript{\cite[][p.~455]{weissMultiagentSystems2013}}
    Supports cooperative, competitive, and mixed environments.

    \textbf{MAPPO} applies PPO (Proximal Policy Optimisation) in the CTDE
    framework with a shared or individual centralised value function.
    Empirically strong on cooperative benchmarks:
    \begin{itemize}
        \item Simple and stable training
        \item On-policy: less sample efficient but more stable than off-policy methods
        \item State-of-the-art on SMAC (StarCraft Multi-Agent Challenge)
    \end{itemize}

\end{v3topic}
\begin{v3topic}{Cooperative vs.\ Competitive vs.\ Mixed}
    MARL environments differ in the structure of the reward
    function:\textsuperscript{\cite[][p.~440]{weissMultiagentSystems2013}}
    \begin{center}
    \begin{tblr}{
        colspec={X[l]X[l]X[l]X[l]},
        hlines, vlines,
        row{1}={font=\bfseries\small, bg=LimeGreen!20},
        row{2-Z}={font=\small},
    }
        Setting & Reward structure & Challenge & Example \\
        Cooperative & Shared team reward & Credit assignment, coordination & Drone rescue swarm \\
        Competitive & Zero-sum (opponents) & Equilibrium finding; opponent modelling & Chess, StarCraft 1v1 \\
        Mixed (general-sum) & Individual rewards & Balance self-interest \& cooperation & Trading agents, autonomous driving \\
        Team competition & Teams share within, compete across & Both challenges combined & AlphaStar, RoboCup \\
    \end{tblr}
\captionof{table}{Key MARL Algorithms: Cooperative vs.\ Competitive vs.\ Mixed.}
\label{tab:ch15-multi-agent-systems-inline-05}
    \end{center}

\end{v3topic}

\subsection{Applications of MARL}

\begin{v3topic}{StarCraft and AlphaStar}
    \textbf{AlphaStar} (Vinyals et al., DeepMind, 2019) achieved Grandmaster
    level in StarCraft II using a multi-agent league of agents that
    train against each other through \emph{self-play and fictitious
    self-play}.\textsuperscript{\cite[][Ch.~18]{weissMultiagentSystems2013}}

    Key techniques:
    \begin{itemize}
        \item \textbf{League training}: diverse pool of agents with different strategies
              prevents collapse to a single strategy
        \item \textbf{Transformer-based} memory over long game trajectories
        \item \textbf{Pointer networks} for selecting units from variable-size sets
        \item Demonstrated that MAS with RL can master complex imperfect-information games
    \end{itemize}

\end{v3topic}
\begin{v3topic}{Multi-Robot Coordination \& Traffic}
    Practical MARL applications beyond games:
    \begin{description}[leftmargin=0.6cm, labelindent=0pt]
        \item[\textbf{Warehouse robots}] MARL agents allocate picking tasks and
              avoid collisions in large flat warehouses (Amazon Robotics)
        \item[\textbf{Multi-UAV mission}] Drone swarms cooperatively search areas,
              track targets, and maintain communication relays
        \item[\textbf{Traffic signal control}] Each intersection is an agent;
              MARL coordinates green phases globally to reduce
              congestion --- outperforms hand-tuned timings by 20--30\%
        \item[\textbf{Energy grid}] Demand-response agents cooperatively balance
              supply from intermittent renewables
    \end{description}

\end{v3topic}

% -------------------------------------------------------
\section{MAS in Supply Chains and Industry}

\subsection{Strategic and Operational Applications}

\begin{v3topic}{Supply Chain Coordination}
    A supply chain is a natural MAS: manufacturers, suppliers, distributors, and
    retailers are autonomous agents with local information and individual
    objectives.\textsuperscript{\cite[][p.~22]{chopraSupplyChainManagement2019}}

    MAS capabilities applied at each planning horizon:
    \begin{description}[leftmargin=0.6cm, labelindent=0pt]
        \item[\textbf{Strategic}] Supplier selection, network design, make-or-buy
              (negotiating agents run auction-based procurement)
        \item[\textbf{Tactical}] Inventory planning, capacity booking, demand
              coordination (blackboard agents share demand signals)
        \item[\textbf{Operational}] Order fulfilment, routing, exception handling
              (reactive agents reschedule in real time on disruption)
    \end{description}
    MAS mitigate the \textbf{bullwhip effect} by sharing point-of-sale data
    up the chain and coordinating replenishment decisions.

\end{v3topic}
\begin{v3topic}{Production Scheduling}
    \textbf{Job-shop scheduling} --- assigning jobs to machines while minimising
    makespan --- is NP-hard. Agent-based approaches decompose the problem:
    \textsuperscript{\cite[][p.~4]{paolucciAgentBasedManufacturing2005}}
    \begin{itemize}
        \item \textbf{Job agents}: bid for machine time through a Contract Net
        \item \textbf{Machine agents}: accept bids, publish schedules
        \item \textbf{Supervisor agent}: resolves conflicts and enforces deadlines
    \end{itemize}
    RL-based scheduling (e.g., using MAPPO or DQN agents per machine) has shown
    competitive performance with classical solvers on dynamic variants where
    job arrival times are
    stochastic.\textsuperscript{\cite{panJobShopSchedulingRL2021}}

\end{v3topic}

\subsection{MAS in Cyber-Physical Systems}

\begin{v3topic}{Agent-Based Manufacturing and JADE}
    \textbf{Cyber-Physical Systems} (CPS) embed software agents in physical
    machines, enabling them to negotiate work allocation
    dynamically.\textsuperscript{\cite[][Ch.~3]{paolucciAgentBasedManufacturing2005}}

    \textbf{JADE} (Java Agent DEvelopment Framework) is the reference FIPA-compliant
    middleware for industrial MAS:
    \begin{itemize}
        \item Provides AMS, DF, and MTS out of the box
        \item Supports distributed deployment across JVM nodes
        \item GUI tool (\texttt{Sniffer}) traces message exchanges
        \item Used in logistics, healthcare records, and smart factory pilots
    \end{itemize}
    Agent containers can run on embedded PLCs and communicate over OPC-UA,
    bridging MAS protocols with industrial automation standards.

\end{v3topic}
\begin{v3topic}{Holonic Manufacturing Systems}
    A \textbf{holon} (Koestler, 1967) is an entity that is simultaneously a
    \emph{whole} (autonomous unit) and a \emph{part} (component of a larger
    system).\textsuperscript{\cite[][p.~102]{paolucciAgentBasedManufacturing2005}}

    \textbf{Holonic Manufacturing Systems} (HMS) model production as a
    hierarchy of holons:
    \begin{description}[leftmargin=0.6cm, labelindent=0pt]
        \item[\textbf{Order holons}] Represent customer orders; negotiate with
              resources to fulfil themselves
        \item[\textbf{Resource holons}] Represent machines or robots; advertise
              their capabilities and accept workloads
        \item[\textbf{Product holons}] Carry manufacturing knowledge of how to produce
              themselves
    \end{description}
    HMS achieves \textbf{plug-and-produce} reconfigurability: adding a new machine
    simply adds a new resource holon to the system.

\end{v3topic}

% -------------------------------------------------------
%
\section{Deep Dive: LLM Agents as Distributed Systems}
\begin{v3topic}{Distinguish Three Agent Traditions}
Classical BDI agents execute explicit beliefs, desires, and intentions.
MARL agents learn policies through interaction.
LLM agents generate plans and tool calls from context.
The three can be combined, but their state, guarantees, and failure modes are
not interchangeable.

\end{v3topic}
\begin{v3topic}{Protocols Need Semantics}
A message schema states syntax; an interaction protocol also defines allowed
states, idempotency, timeouts, authorization, conflict resolution, and recovery.
Use correlation identifiers and explicit ownership so retries do not create
duplicate consequential actions.

\end{v3topic}
\begin{v3topic}{Evaluate Trajectories and Emergence}
Measure task success, steps, tool errors, communication cost, latency,
permission violations, deadlocks, loops, and human interventions.
Test adversarial messages and collusion as well as cooperative examples.
Local pass rates do not prove stable collective behaviour.
\end{v3topic}

\chapterreview{Architect}{Communicate}{Coordinate}{Verify}

\section{Further Reading}

\begin{itemize}
    \item \fullcite{wooldridgeIntroductionMultiAgentSystems2009} --- The standard introductory textbook: agent architectures, BDI, communication, cooperation, game theory, and mechanism design.
    \item \fullcite{weissMultiagentSystems2013} --- Comprehensive MIT Press reference covering agent types, learning, coordination, and MAS applications.
    \item \fullcite{shohamMultiagentSystemsAlgorithmic2009} --- Rigorous treatment of game-theoretic and logical foundations: Nash equilibria, mechanism design, auctions, and social choice.
    \item \fullcite{bellifeministeDevelopingMultiAgentSystems2007} --- Practical guide to building FIPA-compliant MAS with JADE: architecture, interaction protocols, and deployment.
    \item \fullcite{bordiniMultiAgentProgramming2009} --- Languages, tools, and applications for programming MAS: AgentSpeak, Jason, 2APL, and case studies.
    \item \fullcite{paolucciAgentBasedManufacturing2005} --- Application-focused: agent-based manufacturing, holonic systems, and CPS integration.
    \item \fullcite{lanhamAIAgentsAction2025} --- Modern perspective on building production multi-agent systems with LLM-based frameworks (LangChain, AutoGen, CrewAI).
    \item \fullcite{talukdarASYNCHRONOUSTEAMSCOOPERATION} --- Original paper on Asynchronous Teams: cooperation schemes for autonomous agents solving combinatorial optimisation.
    \item \fullcite{suttonReinforcementLearningIntroduction2018} --- Foundational RL textbook; Chapter 17 extends to multi-agent settings.
    \item \fullcite{chopraSupplyChainManagement2019} --- Supply chain management strategy and operations; context for agent-based coordination applications.
\end{itemize}
